\documentclass[11pt]{article}

\usepackage[margin=1in]{geometry}
\usepackage[T1]{fontenc}
\usepackage[utf8]{inputenc}
\usepackage{lmodern}
\usepackage{microtype}
\usepackage{amsmath,amssymb,amsthm}
\usepackage{mathtools}
\usepackage{xcolor}
\usepackage{hyperref}
\usepackage{booktabs}
\usepackage{enumitem}
\usepackage{array}
\usepackage{longtable}

\hypersetup{
  colorlinks=true,
  linkcolor=blue!50!black,
  citecolor=blue!50!black,
  urlcolor=blue!50!black,
  pdftitle={MeTTa Pure/Core Primer},
  pdfauthor={OpenAI GPT-5.2 Pro},
  pdfsubject={Semantic waist, declaration-aware extensions, and proof/execution split for MeTTa Pure/Core},
  pdfkeywords={MeTTa, PureKernel, dependent type theory, declaration semantics, recursors, MM2, MORK}
}

\newcommand{\MeTTa}{\textup{MeTTa}}
\newcommand{\MeTTaIL}{\textup{MeTTaIL}}
\newcommand{\MORK}{\textup{MORK}}
\newcommand{\MMTwo}{\textup{MM2}}
\newcommand{\code}[1]{\texttt{#1}}
\newcommand{\defeq}{\ensuremath{\mathrel{\equiv}}}
\newcommand{\TypeZero}{\ensuremath{\mathsf{Type0}}}
\newcommand{\TypeOne}{\ensuremath{\mathsf{Type1}}}
\newcommand{\Pure}{\textsf{PureKernel}}
\newcommand{\Core}{\textsf{ElaboratedCore}}
\newcommand{\cdev}{\ensuremath{\mathsf{cdev}}}
\newcommand{\deltared}{\ensuremath{\delta}}
\newcommand{\iotared}{\ensuremath{\iota}}
\newcommand{\Nat}{\ensuremath{\mathsf{Nat}}}
\newcommand{\Bool}{\ensuremath{\mathsf{Bool}}}
\newcommand{\UnitTy}{\ensuremath{\mathsf{Unit}}}
\newcommand{\ListTy}{\ensuremath{\mathsf{List}}}
\newcommand{\VecTy}{\ensuremath{\mathsf{Vec}}}

\title{A Primer on \MeTTa{} Pure/Core\\
\large The Semantic Waist, Declaration-Aware Families, and the Proof/Execution Split}
\author{AI-assisted conceptual note for continuity across development sessions}
\date{2026-03-13}

\begin{document}
\maketitle

\begin{abstract}
This note consolidates the main architectural and mathematical ideas that emerged
while refining the \MeTTa{} Pure/Core development. Its purpose is not to record
commit history. Instead, it isolates the durable conceptual picture that should
continue to guide future work.

The central thesis is that the project should converge to a \emph{small trusted
semantic waist} on the proof side, with a separate runtime execution seam on the
operational side, both fed from a shared elaborated core. The proof-side waist
should remain generic and theoremic; ordinary families, recursors, and
structural recursion should arise from declaration-aware machinery above that
waist rather than from family-specific syntax hacks. Prototype evaluators and
surface conveniences may exist, but they must be downstream of the trusted
kernel rather than silently shaping it.

The note explains the clean architecture, the main lessons learned from drift and
recovery, the declaration-aware route to ordinary families, the correct shape of
generated dependent recursors, the role of positivity and substitution, the
separation between $\delta$-unfolding and $\iota$-reduction, and the current
status of the connection to \MeTTaIL{} and the \MMTwo{}/\MORK{} runtime seam.
It is intended to serve as a didactic primer for future AI-assisted development
sessions.
\end{abstract}

\tableofcontents

\section{Purpose and Scope}

The goal of this note is to make the Pure/Core project legible in a way that is
stable across individual coding sessions. The emphasis is on \emph{conceptual
architecture}, not on temporary implementation detail.

The note has four jobs:
\begin{enumerate}[leftmargin=2em]
  \item identify the mathematically clean center of the project;
  \item explain which extensions belong \emph{above} that center rather than
  \emph{inside} it;
  \item record the most important failure modes that must not be repeated;
  \item provide a compact primer that future development sessions can use as a
  normative reference.
\end{enumerate}

The intended reader is someone who may be new to the repository but needs to
understand the design constraints immediately and accurately.

\section{The Architectural Thesis}

The working architectural thesis can be stated in one sentence:

\begin{quote}
Surface \MeTTa{} should elaborate to a shared middle layer, and that shared
middle layer should support two distinct but connected targets: a proof-oriented
semantic target and a runtime execution target.
\end{quote}

The three main layers are therefore:
\begin{enumerate}[leftmargin=2em]
  \item \textbf{Surface \MeTTa.} The user-facing language, notation, parser,
  syntax sugar, and future ergonomic layer.
  \item \textbf{Elaborated Core.} A shared, explicit intermediate layer used to
  mediate between syntax and semantics.
  \item \textbf{Two honest targets.}
  \begin{enumerate}[label=\alph*.,leftmargin=2em]
    \item a proof-side target, centered on a small trusted dependent kernel;
    \item an execution-side target, centered on a real operational seam such as
    the \MMTwo{}/\MORK{} runtime stack.
  \end{enumerate}
\end{enumerate}

This architecture matters because it prevents two opposite confusions:
\begin{enumerate}[leftmargin=2em]
  \item \emph{runtime reduction is not the proof kernel};
  \item \emph{proof semantics and runtime execution are not permanently
  unrelated silos either}.
\end{enumerate}

The shared elaborated core is the place where those worlds can meet honestly.

\section{The Semantic Waist}

\subsection{What the waist is}

A \emph{semantic waist} is the smallest trusted layer that is still rich enough
to support the key theoremic services of the proof side. In the present design,
that waist is a small dependent kernel with:
\begin{itemize}[leftmargin=2em]
  \item universes (currently the tiny \TypeZero{}/\TypeOne{} story),
  \item dependent function types ($\Pi$),
  \item dependent pair types ($\Sigma$),
  \item identity types,
  \item ordinary variables and binders,
  \item global constants,
  \item a small, explicit reduction relation,
  \item typing, confluence, subject reduction, and definitional equality.
\end{itemize}

The crucial point is that this waist should be \emph{generic}. It should not be
polluted by family-specific constructors such as hardcoded \Nat{}, \Bool{}, or
recursor nodes if the intended final architecture is a general family mechanism.

\subsection{Why the waist must stay small}

A small waist has four advantages.
\begin{enumerate}[leftmargin=2em]
  \item \textbf{Metatheory stays local.} Confluence, subject reduction,
  algorithmic typing, and conversion are harder to trust when the syntax grows
  opportunistically.
  \item \textbf{Extensions stay honest.} New user-visible features must justify
  themselves as elaboration or declaration mechanisms above the waist rather
  than by mutating the kernel whenever a demo is desired.
  \item \textbf{The trust boundary is explicit.} One can say exactly what is
  trusted and what is merely derived.
  \item \textbf{Future automation becomes safer.} AI agents are much less
  likely to introduce architectural drift when the kernel surface is visibly
  narrow and well guarded.
\end{enumerate}

\subsection{What the waist is not}

The semantic waist is \emph{not}:
\begin{itemize}[leftmargin=2em]
  \item a parser;
  \item a convenience evaluator;
  \item a proof explorer UI;
  \item a staging area for temporary built-ins;
  \item a promise that every future user-level construct already has kernel
  computation rules.
\end{itemize}

Those can all exist, but none of them should define the trusted core.

\section{The Proof-Side Reference Must Be Theoremic, Not Evaluator-Driven}

One of the most important lessons is that the proof-side \emph{reference} should
be the checked theoremic boundary, not a prototype evaluator.

The conceptual center of the proof side should be:
\begin{enumerate}[leftmargin=2em]
  \item typed checking into the trusted kernel;
  \item canonicalization by a theoremically sanctioned notion such as complete
  development (written here as $\cdev$);
  \item checked bridges between syntax, kernel terms, and shared artifacts.
\end{enumerate}

A small evaluator can still be useful, but only if it is treated as a derived
explorer. The key rule is:

\begin{quote}
Prototype evaluation may \emph{test} the kernel, but it may not \emph{shape}
the kernel.
\end{quote}

This distinction matters because evaluator-driven development tends to create
exactly the wrong incentives. It rewards short-term demonstrations and can make
non-final shortcuts look like conceptual progress. The clean alternative is to
let theoremic checking and canonicalization remain authoritative, while any
executable evaluator is explicitly secondary.

\section{Legacy Ambient Presentations versus the Trusted Kernel}

It is perfectly reasonable for a project like this to contain more than one
presentation of ``Pure'':
\begin{itemize}[leftmargin=2em]
  \item a legacy or ambient presentation over a broader host language;
  \item a smaller, more explicit kernel-level presentation used as the trusted
  center.
\end{itemize}

The danger is not the existence of both. The danger is forgetting which one is
trusted.

The decisive lesson from the recent confluence repair is the following:

\begin{quote}
A theorem, its reduction relation, and its fragment assumptions must all live on
exactly the same fragment.
\end{quote}

If confluence is only proved for a fragment of terms, then the conversion
relation used by downstream typing must be fragment-safe as well. One must not
silently claim Church--Rosser for a broader ambient relation and then sneak in
fragment facts only at the call site. The correct move is to tighten the
relation itself or tighten the theorem boundary, not to fabricate missing
hypotheses.

This is more than a technicality. It is a design law:
\begin{enumerate}[leftmargin=2em]
  \item fragment membership should be an upstream concept,
  \item binder and purity facts should be upstream utilities,
  \item confluence should be downstream of those facts,
  \item subject reduction should be downstream of confluence,
  \item no proof should quietly range over a larger syntax than its lemmas can
  actually support.
\end{enumerate}

\section{Declaration-Aware Extensions above the Waist}

\subsection{Why declarations are the right extension point}

Once the small kernel includes \emph{global constants}, the clean next step is
not to add more syntax constructors. It is to add a declaration-aware layer
above the waist.

Conceptually, this means introducing:
\begin{itemize}[leftmargin=2em]
  \item an environment of declarations,
  \item a declaration-aware typing judgment,
  \item a declaration-aware reduction relation,
  \item well-formedness conditions on the environment.
\end{itemize}

This allows one to treat families, constructors, and recursors as ordinary named
declarations rather than new kernel AST nodes.

\subsection{The current layering idea}

A principled layering for declaration-aware work is:
\begin{enumerate}[leftmargin=2em]
  \item \textbf{Minimal operational environment well-formedness.}
  Enough to justify lookup and basic $\deltared$-unfolding soundly.
  \item \textbf{Ordered signatures.}
  Enough to talk about declaration order, shadowing discipline, and prefix
  reasoning.
  \item \textbf{Prefix discipline.}
  Enough to say that declaration types and unfolding values only mention names
  from earlier entries.
  \item \textbf{Full family-formation theory.}
  A later layer that justifies families, constructors, generated recursors, and
  equation schemas as a genuine semantic mechanism rather than a list of pilots.
\end{enumerate}

This is a good architecture because it does not pretend the first operational
invariant is already the whole declaration theory. The intermediate layers make
order semantically meaningful before the full family machinery is present.

\subsection{Why prefix discipline matters}

A pure ``global final environment'' view is too weak. It can hide circularity,
self-reference, and order-dependence. Prefix discipline repairs this by making
left-to-right growth real:
\begin{itemize}[leftmargin=2em]
  \item fresh names against the prefix,
  \item values typecheck against the prefix,
  \item declaration types mention only earlier names,
  \item declaration values mention only earlier names,
  \item no self-unfolding at the point of introduction.
\end{itemize}

This is not yet the final family theory. It is an honest bridge toward it.

\subsection{What still needs to become first-class}

The next theoremic bundling should introduce a stronger environment invariant,
for example a conceptually named declaration-environment well-formedness record,
rather than leaving individual obligations scattered.

Such a record should eventually bundle at least:
\begin{itemize}[leftmargin=2em]
  \item well-formed declared types,
  \item well-typed declared values,
  \item closure/scoping requirements,
  \item a stratification or acyclicity discipline if unfolding becomes more
  central to conversion.
\end{itemize}

\section{Ordinary Families Must Be Declarations, Not Kernel Hacks}

\subsection{The main principle}

Ordinary inductive families such as \UnitTy{}, \Bool{}, \Nat{}, and later
indexed families such as $\VecTy$, should be \emph{instances of a general
family-declaration mechanism}. They should not be introduced by mutating the
trusted syntax of the small kernel whenever a new example is desired.

This principle has a direct corollary:

\begin{quote}
If a feature is meant to come from a later general mechanism, it should not be
hardcoded now as a family-specific shortcut in the trusted lane.
\end{quote}

\subsection{What a genuine family declaration must contain}

A real ordinary-family declaration layer should be able to express, at minimum:
\begin{itemize}[leftmargin=2em]
  \item the family name,
  \item a \emph{parameter telescope},
  \item an \emph{index telescope},
  \item the resulting sort,
  \item constructors with full \emph{typed telescopes},
  \item a \emph{generated recursor} specification,
  \item \emph{equation schemas} for generated computation,
  \item \emph{positivity/admissibility} evidence.
\end{itemize}

Several important negative conclusions follow immediately:
\begin{enumerate}[leftmargin=2em]
  \item a list of closed types is not enough when dependency matters;
  \item constructor arity is not constructor typing;
  \item a closed left-hand side/right-hand side pair is not yet a generic
  equation schema under binders;
  \item boolean flags are not acceptable substitutes for proof obligations in
  the trusted lane.
\end{enumerate}

\subsection{Dependent telescopes are mandatory}

If telescopes are represented by closed terms only, then later binders cannot
depend on earlier ones. That is already too weak for ordinary family work.

So the family record should use a genuinely binder-aware telescope rather than a
flat list of closed types. The family constant type should then be derived from
that telescope by iterated $\Pi$-formation, rather than stored redundantly and
risking divergence between two representations.

\section{Dependent Recursors Are the Canonical Target}

\subsection{The canonical shape}

For a non-indexed family $F$ with parameters $\Gamma_p$ and constructors in
order $c_1,\dots,c_k$, the canonical generated recursor should be the dependent
eliminator:
\[
\Pi \Gamma_p.
\Pi P : F\,\Gamma_p \to \TypeZero.
\Pi \mathit{case}_1 : \mathit{CaseTy}_1.
\cdots
\Pi \mathit{case}_k : \mathit{CaseTy}_k.
\Pi x : F\,\Gamma_p.
P\,x.
\]

For indexed families, the parameters still come \emph{before} the motive, while
the indices are bound \emph{inside} the motive and major premise. In other
words, the right rule is:
\begin{quote}
parameters outside the motive; indices inside the motive.
\end{quote}

Case arms should appear in constructor declaration order. This keeps the
relationship between declarations, generated types, and generated computation
rules stable and reviewable.

\subsection{The \Nat{} example}

The canonical dependent eliminator for naturals is:
\[
\Nat.\mathsf{rec} :
\Pi P : \Nat \to \TypeZero.
P\,(\Nat.\mathsf{zero}) \to
\bigl(\Pi n : \Nat.\; P\,n \to P\,(\Nat.\mathsf{succ}\;n)\bigr) \to
\Pi n : \Nat.\; P\,n.
\]

A non-dependent programming recursor,
\[
\Pi R : \TypeZero.\; R \to (\Nat \to R \to R) \to \Nat \to R,
\]
may be a useful \emph{derived interface}, but it should not be the primary
generator target. The dependent eliminator is the canonical source of truth.

\subsection{The \ListTy{} example}

For a future list family $\ListTy : \Pi A : \TypeZero.\; \TypeZero$, the
canonical dependent eliminator is:
\[
\ListTy.\mathsf{rec} :
\Pi A : \TypeZero.
\Pi P : \ListTy\,A \to \TypeZero.
P\,(\ListTy.\mathsf{nil}\;A) \to
\bigl(\Pi x : A.\; \Pi xs : \ListTy\,A.\; P\,xs \to P\,(\ListTy.\mathsf{cons}\;A\;x\;xs)\bigr) \to
\Pi xs : \ListTy\,A.\; P\,xs.
\]

This is the right conceptual template for future generated recursors.

\section{Positivity and Indexed Families}

Strict positivity should be checked over the \emph{full constructor type}, not
just over the parametric part. Once indices are present, recursive occurrences
can hide in index-bearing targets and argument positions that a parameter-only
check would miss.

However, positivity and target-shape checking should be separated into two
orthogonal constraints:
\begin{enumerate}[leftmargin=2em]
  \item \textbf{Uniform parameters.} The parameter prefix of every recursive
  family occurrence must be exactly the declared parameters, in the expected
  positions.
  \item \textbf{Flexible indices.} Index arguments may be arbitrary well-scoped
  terms over the constructor context.
\end{enumerate}

This matters for indexed families such as
\[
\VecTy : \Pi A : \TypeZero.\; \Nat \to \TypeZero,
\]
where a target like $\VecTy\;A\;(\Nat.\mathsf{succ}\;n)$ must be allowed.

A clean implementation strategy is therefore to decompose applications into a
head and an argument spine, check the family head, split the arguments into
parameter and index portions, enforce exact matching on the parameter prefix,
and permit arbitrary scoped terms for the indices.

\section{Substitution and Generated Computation}

\subsection{Parallel substitution is the right semantic primitive}

A useful conceptual simplification is this: if the kernel already has parallel
substitution as its basic operation, then generated recursor reduction does not
require a new foundation. It requires a new \emph{wrapper}.

In other words, a simultaneous substitution API is semantically sufficient.
What is usually missing is only an ergonomic interface for substituting several
constructor arguments and induction hypotheses into a case body in the right
de Bruijn order.

A thin convenience operation such as \code{substVec} or \code{instMany} is a
good idea. But it should be understood as a thin layer over the existing
parallel substitution theory rather than as a new primitive.

\subsection{The correct design rule}

The important design rule is:
\begin{quote}
Do not rebuild the semantics around repeated one-variable substitution if a
clean simultaneous substitution is already available.
\end{quote}

Repeated one-variable substitution is much easier to get wrong, especially under
binders and when constructor arguments and induction hypotheses interleave.

\section{$\delta$-Unfolding versus $\iota$-Reduction}

A declaration-aware kernel needs to distinguish two different mechanisms:
\begin{enumerate}[leftmargin=2em]
  \item \textbf{$\deltared$-unfolding.} A named constant unfolds to its stored
  value by declaration lookup.
  \item \textbf{$\iotared$-reduction.} A generated recursor computes on a
  saturated application of the recursor to a constructor.
\end{enumerate}

These are not the same thing, and they should not be modeled by the same data
field.

\subsection{What should stay in a declaration specification}

A declaration specification should remain focused on the ordinary environment
layer:
\begin{itemize}[leftmargin=2em]
  \item name,
  \item type,
  \item optional value for ordinary $\deltared$-unfolding,
  \item well-formedness conditions relative to the current signature.
\end{itemize}

\subsection{What should live above it}

Generated recursor computation should live in a \emph{separate} semantic layer,
for example a generated-recursor rule bundle or family-level computation record.
That higher layer should package:
\begin{itemize}[leftmargin=2em]
  \item the generated recursor declaration,
  \item the generated $\iotared$-rules,
  \item typing proofs for those rules,
  \item coverage and arity discipline,
  \item whatever later confluence or conversion facts are needed.
\end{itemize}

This separation keeps the ordered-signature story clean. It also prevents
constant lookup from being muddled with eliminator computation.

\section{The Runtime Frontier and the Meaning of ``a Strain of \MeTTa{}''}

The Pure/Core project is not aiming to replace the runtime semantics with proof
semantics. The goal is subtler.

The proof-side target should remain a mathematically isolated dependent kernel.
The execution-side target should remain a real operational seam with genuine
runtime behavior, currently exemplified by the \MMTwo{}/\MORK{} line.

The interesting question is therefore not ``is the proof kernel literally the
runtime?'' but rather:
\begin{quote}
What is the cleanest elaborated-core fragment that can lower honestly to both
sides?
\end{quote}

This reframes the relation between Pure and execution in a productive way.

\subsection{The honest current status}

The current honest position is:
\begin{itemize}[leftmargin=2em]
  \item the proof side already has a trusted checking and canonicalization
  center;
  \item the runtime side already has a real execution seam;
  \item there is not yet a theoremically established direct Pure-to-runtime
  execution bridge for the full proof fragment;
  \item the next mathematically meaningful connection is an \emph{intentional
  dual-target fragment} of elaborated core.
\end{itemize}

This is a healthy status, not a deficiency. It keeps the architecture honest.

\section{Lessons from Drift and Recovery}

Several hard lessons emerged from the development process. They are worth
recording explicitly.

\subsection{Progress can be fake if the trust boundary moves}

A project can appear to make fast progress while actually losing architectural
integrity. Typical examples include:
\begin{itemize}[leftmargin=2em]
  \item adding family-specific kernel constructors because a demonstration seems
  easier that way;
  \item promoting a convenience evaluator to reference status;
  \item widening a theorem statement beyond the fragment actually proved;
  \item leaving stale examples and docs in place after a semantic reset;
  \item putting pilot modules on the same import surface as the trusted core.
\end{itemize}

The appearance of momentum is not enough. The relevant question is always:
\begin{quote}
Did the new capability arise from the intended final mechanism, or from a
shortcut that would have to be deleted later?
\end{quote}

\subsection{The right response to a bad shortcut is usually retreat, not polish}

When a trusted layer has absorbed a shortcut that is not part of the intended
final design, the correct move is usually to back it out rather than to make it
nicer. Cosmetic polish does not convert a non-final mechanism into a final one.

\subsection{A good proof failure is valuable}

The confluence break in the ambient legacy layer was a useful event because it
revealed that the theorem boundary was overclaimed. Proof failures that expose
mismatched abstraction boundaries are healthy. They should trigger conceptual
repair, not theorem-patching theater.

\section{Development Discipline for Future AI-Assisted Sessions}

The project needs explicit design rules that future sessions can enforce.

\subsection{Core rules}

\begin{enumerate}[leftmargin=2em,label=\textbf{R\arabic*.}]
  \item \textbf{No temporary trusted constructors.}
  Never add a trusted kernel constructor, eliminator, typing rule, or reduction
  rule unless it belongs to the intended final architecture.

  \item \textbf{No evaluator-driven semantics.}
  Prototype evaluators may exist, but they may not determine the trusted kernel
  shape.

  \item \textbf{No family-specific shortcuts when a general mechanism is the
  goal.}
  If a feature is supposed to arise from a generic family declaration system,
  do not hardcode a family-specific version in the trusted lane.

  \item \textbf{Match theorem and fragment.}
  The syntax fragment, conversion relation, and theorem statement must range
  over the same object language.

  \item \textbf{Keep trust labels explicit.}
  Every new component should be classifiable as trusted kernel, declaration
  extension, pilot/specification, or experimental tooling.

  \item \textbf{Synchronize code, docs, and examples.}
  Any semantic change must update examples and documentation in the same pass.

  \item \textbf{Prefer proof objects to booleans in the trusted lane.}
  A boolean flag may be acceptable during exploration, but the trusted theory
  should expose propositions, structures, and theorems rather than hope encoded
  as data.
\end{enumerate}

\subsection{Questions every future change should answer}

Before accepting a kernel-adjacent change, one should be able to answer:
\begin{enumerate}[leftmargin=2em]
  \item Does this belong inside the small waist, or above it?
  \item If it touches syntax, why can constants plus declarations not express it
  instead?
  \item Is the new semantics theoremic, or is it being inferred from a tool?
  \item Are the fragment assumptions of all downstream theorems still honest?
  \item Does this make future generalization easier, or does it create more
  bespoke cases?
\end{enumerate}

\section{A Robustness Checklist}

The Pure/Core project can be considered to have reached a robust kernel stage
when the following become true.

\begin{enumerate}[leftmargin=2em]
  \item \textbf{New families are declaration instances, not syntax changes.}
  Adding \Nat{}, \Bool{}, \ListTy{}, or an indexed family should no longer
  require edits to the trusted kernel AST.

  \item \textbf{Recursors are generated, not hand-added.}
  The dependent recursor for a family should come from a generic generation
  procedure justified by the family declaration mechanism.

  \item \textbf{The declaration preservation stack is complete.}
  In particular: one-step $\deltared$ preservation, one-step declaration-aware
  preservation, and star preservation.

  \item \textbf{The proof-side reference is stable.}
  The checked kernel plus canonicalization remains the reference; evaluators and
  UI layers remain downstream.

  \item \textbf{The bridge to execution is intentional rather than accidental.}
  There is at least one explicitly identified elaborated-core fragment that
  lowers honestly to both the proof side and the runtime side.
\end{enumerate}

\section{A Suggested Near-Term Roadmap}

The following sequence is conceptually conservative and mathematically aligned
with the architecture.

\begin{enumerate}[leftmargin=2em]
  \item \textbf{Keep the small kernel frozen.}
  No family-specific syntax growth.

  \item \textbf{Make trust boundaries visible.}
  Separate tiny trusted-kernel umbrellas from declaration-extension umbrellas
  and pilot/experimental umbrellas.

  \item \textbf{Bundle declaration-environment well-formedness more tightly.}
  Turn the current scattered obligations into an explicit semantic record.

  \item \textbf{Complete the declaration preservation tower.}
  Prove one-step $\deltared$ preservation, then general declaration-aware
  one-step preservation, then star preservation.

  \item \textbf{Introduce a binder-aware generic family declaration record.}
  Use dependent telescopes, derived family types, typed constructors, generated
  dependent recursors, equation schemas, and proof-valued admissibility fields.

  \item \textbf{Realize \UnitTy{} as the first real instance.}
  It is the simplest family and the best proving ground for the generic record.

  \item \textbf{Add generated $\iotared$ semantics above the signature layer.}
  Keep ordinary declaration lookup and generated recursor computation distinct.

  \item \textbf{Then add \Bool{}, then \Nat{}.}
  Each should become an instance of the same mechanism, not a special case.

  \item \textbf{Only afterwards revisit stronger canonicalization and the proof/
  runtime frontier.}
  By then, the new family layer will be stable enough to connect outward.
\end{enumerate}

\section{Glossary of Key Terms}

\begin{description}[leftmargin=2em,style=nextline]
  \item[Semantic waist] The smallest trusted proof-side layer rich enough to
  support checking, canonicalization, and later extensions.

  \item[Elaborated core] The explicit shared middle layer between surface
  \MeTTa{} and the proof/runtime targets.

  \item[Canonicalization / $\cdev$] The theoremically authoritative proof-side
  reduction or normalization anchor, distinct from any convenience evaluator.

  \item[$\deltared$-unfolding] Constant unfolding by declaration lookup.

  \item[$\iotared$-reduction] Generated eliminator computation on a recursor
  applied to a constructor.

  \item[Declaration-aware kernel] The small kernel together with declaration
  environments, declaration-aware typing, and declaration-aware reduction.

  \item[Prefix discipline] The requirement that declaration types and values may
  only mention earlier names, so signature order has semantic force.

  \item[Ordinary family] A strictly positive inductive family, possibly with
  parameters and later with indices.

  \item[Dependent recursor] The canonical eliminator whose motive may depend on
  the major premise (and, for indexed families, on the indices).

  \item[Dual-target fragment] A deliberately chosen fragment of elaborated core
  that lowers honestly to both the proof side and the runtime execution side.
\end{description}

\section{Conclusion}

The main conceptual picture is now clear.

The Pure/Core project should be organized around a small generic trusted kernel,
a declaration-aware extension above that kernel, a future general family
mechanism that treats families as declarations rather than syntax hacks, and a
separate runtime execution seam connected through a shared elaborated core.

The most important design lesson is that correctness and architectural honesty
must outrank short-term demonstrations. A project like this becomes durable only
when each layer is justified by the intended final mechanism, not by temporary
convenience.

If future development keeps the semantic waist small, keeps declarations ordered
and meaningful, generates families and recursors from general records, and
connects to execution through explicit dual-target fragments rather than wishful
identification, then the project has a credible path toward a robust and
mathematically well-isolated version of \MeTTa{} Pure/Core.

\end{document}
